\documentclass[a4paper,landscape,8pt]{extarticle}
\usepackage[landscape,margin=0.7cm,top=0.5cm,bottom=0.7cm]{geometry}
\usepackage[utf8]{inputenc}
\usepackage[T1]{fontenc}
\usepackage[ngerman]{babel}
\usepackage{helvet}
\renewcommand{\familydefault}{\sfdefault}
\usepackage{xcolor}
\usepackage{tikz}
\usetikzlibrary{positioning, shapes.geometric}
\usepackage{enumitem}
\usepackage{multicol}
\usepackage{array}
\usepackage{fancyhdr}
\usepackage{amssymb}

% Farben
\definecolor{lockgreen}{HTML}{34C759}
\definecolor{keyblue}{HTML}{007AFF}
\definecolor{shieldgray}{HTML}{8E8E93}
\definecolor{warnred}{HTML}{C62828}
\definecolor{weakorange}{HTML}{FF9500}
\definecolor{lightgray}{HTML}{F5F5F5}
\definecolor{bordergray}{HTML}{BBBBBB}

\pagestyle{fancy}
\fancyhf{}
\renewcommand{\headrulewidth}{0pt}
\fancyfoot[C]{\footnotesize Seite \thepage{} von 3}

\setlength{\parindent}{0pt}
\setlength{\parskip}{2pt}

% Kompakte Listen
\setlist{nosep, leftmargin=1em, itemsep=0pt, topsep=1pt}

\begin{document}

% ============================================================================
% SEITE 1: INFORMATIONSBLATT
% ============================================================================

\noindent{\Large\bfseries Passwörter: Dein digitaler Schlüsselbund} \hfill
{\footnotesize\textbf{Lernziele:} Sichere Passwörter | Schlüsselbund nutzen | Passwörter verwalten}

\vspace{1mm}
\hrule
\vspace{1mm}

% Drei-Spalten-Layout
\begin{minipage}[t]{0.32\linewidth}
\begin{center}
\colorbox{warnred!10}{%
\begin{minipage}{0.95\linewidth}
\vspace{2mm}
\begin{center}
{\Large\bfseries\textcolor{warnred}{Warum ist das wichtig?}}\\[1pt]
{\footnotesize\itshape Passwörter schützen dein Leben}
\end{center}
\vspace{2mm}
\end{minipage}%
}
\end{center}

\vspace{1mm}

\textbf{Was schützen Passwörter?}
\begin{itemize}
\item Deine \textbf{E-Mails} (private Nachrichten)
\item Dein \textbf{Bankkonto} (Geld!)
\item Deine \textbf{Fotos} (Erinnerungen)
\item Deine \textbf{Identität} (Missbrauch verhindern)
\end{itemize}

\textbf{Analogie:} Passwörter sind wie \textbf{Schlüssel} -- jeder Schlüssel öffnet eine andere Tür.

\vspace{2mm}

\textbf{Was passiert bei schlechten Passwörtern?}
\begin{itemize}
\item Fremde lesen deine E-Mails
\item Kriminelle kaufen mit deinem Geld ein
\item Jemand gibt sich als du aus
\item Deine Daten werden verkauft
\end{itemize}

\vspace{2mm}

\begin{center}
\fcolorbox{warnred}{warnred!8}{%
\begin{minipage}{0.92\linewidth}
\vspace{2mm}
\textbf{Wichtigste Regel:}\\[2pt]
\textbf{Niemals} das gleiche Passwort für mehrere Konten verwenden!
\vspace{2mm}
\end{minipage}}
\end{center}

\end{minipage}%
\hfill%
\begin{minipage}[t]{0.32\linewidth}
\begin{center}
\colorbox{lockgreen!10}{%
\begin{minipage}{0.95\linewidth}
\vspace{2mm}
\begin{center}
{\Large\bfseries\textcolor{lockgreen}{Sichere Passwörter}}\\[1pt]
{\footnotesize\itshape Was macht ein gutes Passwort aus?}
\end{center}
\vspace{2mm}
\end{minipage}%
}
\end{center}

\vspace{1mm}

\textbf{Ein sicheres Passwort hat:}
\begin{itemize}
\item Mindestens \textbf{12 Zeichen} (besser: 16+)
\item \textbf{Großbuchstaben} (A, B, C...)
\item \textbf{Kleinbuchstaben} (a, b, c...)
\item \textbf{Zahlen} (1, 2, 3...)
\item \textbf{Sonderzeichen} (!, @, \#, \$...)
\end{itemize}

\textbf{Gute Methode -- der Satz-Trick:}

Denke dir einen Satz aus und nimm die Anfangsbuchstaben:

\glqq \textit{Meine Oma backt 3 Kuchen am Sonntag!}\grqq{}

$\rightarrow$ \texttt{MOb3KaS!}

Noch besser mit mehr Zeichen:

\glqq \textit{Ich trinke jeden Morgen 2 Tassen Kaffee!}\grqq{}

$\rightarrow$ \texttt{ItjM2TK!}

\vspace{2mm}

\textbf{Oder:} Lass dir vom iPhone ein \textbf{sicheres Passwort vorschlagen} -- es merkt sich das für dich!

\end{minipage}%
\hfill%
\begin{minipage}[t]{0.32\linewidth}
\begin{center}
\colorbox{weakorange!10}{%
\begin{minipage}{0.95\linewidth}
\vspace{2mm}
\begin{center}
{\Large\bfseries\textcolor{weakorange}{Häufige Fehler}}\\[1pt]
{\footnotesize\itshape Das solltest du vermeiden}
\end{center}
\vspace{2mm}
\end{minipage}%
}
\end{center}

\vspace{1mm}

\textbf{Schlechte Passwörter:}
\begin{itemize}
\item \texttt{123456} -- zu einfach
\item \texttt{passwort} -- zu offensichtlich
\item \texttt{Geburtsdatum} -- leicht zu erraten
\item \texttt{Name des Hundes} -- findet man heraus
\item \texttt{qwertz} -- Tastaturmuster
\end{itemize}

\textbf{Diese Fehler vermeiden:}
\begin{itemize}
\item[$\times$] Gleiches Passwort überall
\item[$\times$] Passwort auf Zettel am Monitor
\item[$\times$] Passwort per E-Mail schicken
\item[$\times$] Passwort am Telefon sagen
\item[$\times$] Zu kurze Passwörter (unter 8 Zeichen)
\end{itemize}

\vspace{2mm}

\textbf{Passwort-Stärke:}

\renewcommand{\arraystretch}{1.3}
\begin{tabular}{@{}p{2cm}p{4cm}@{}}
\textcolor{warnred}{Schwach} & \texttt{hund123} \\
\textcolor{weakorange}{Mittel} & \texttt{Hund\#123} \\
\textcolor{lockgreen}{Stark} & \texttt{M3inHund!Li3btK\#s3} \\
\end{tabular}
\renewcommand{\arraystretch}{1}

\end{minipage}

\vspace{1mm}

% Zusammenfassung
\begin{center}
\fcolorbox{bordergray}{lightgray}{%
\begin{minipage}{0.98\linewidth}
\centering
\vspace{1.5mm}
\textbf{Merke:} Jedes Konto braucht ein \textbf{eigenes, starkes Passwort}. Nutze den \textbf{Satz-Trick} oder lass dir Passwörter \textbf{vorschlagen und speichern}.
\vspace{1.5mm}
\end{minipage}%
}
\end{center}

\newpage

% ============================================================================
% SEITE 2: SCHLÜSSELBUND & PASSWORT-MANAGER
% ============================================================================

\begin{center}
{\LARGE\bfseries Der iCloud-Schlüsselbund \& Passwörter verwalten}\\[3pt]
{\normalsize Apple merkt sich deine Passwörter -- sicher und bequem}
\end{center}

\vspace{1mm}
\hrule
\vspace{2mm}

\begin{multicols}{2}

\section*{\textcolor{keyblue}{Der iCloud-Schlüsselbund}}

\textbf{Was ist das?} Apples eingebauter \textbf{Passwort-Tresor} -- speichert alle deine Passwörter sicher.

\textbf{Analogie:} Wie ein \textbf{Schlüsselbund} in einem Safe -- alle Schlüssel an einem sicheren Ort.

\vspace{2mm}

\textbf{Vorteile:}
\begin{itemize}
\item Passwörter werden \textbf{verschlüsselt} gespeichert
\item \textbf{Automatisch ausgefüllt} auf Webseiten und in Apps
\item Auf \textbf{allen Apple-Geräten} verfügbar (iPhone, iPad, Mac)
\item Du musst dir nur \textbf{ein Passwort} merken (Geräte-Code)
\end{itemize}

\textbf{Schlüsselbund aktivieren:}
\begin{enumerate}
\item Einstellungen → [Dein Name] → iCloud
\item \textbf{Passwörter \& Schlüsselbund} antippen
\item Einschalten (grüner Schalter)
\end{enumerate}

\vspace{3mm}

\section*{\textcolor{lockgreen}{Passwörter vorschlagen lassen}}

\textbf{Wenn du ein neues Konto erstellst:}
\begin{enumerate}
\item Das iPhone schlägt automatisch ein \textbf{sicheres Passwort} vor
\item Tippe auf \glqq Starkes Passwort verwenden\grqq{}
\item Das Passwort wird \textbf{automatisch gespeichert}
\item Du musst es dir \textbf{nicht merken}!
\end{enumerate}

\textbf{Beim nächsten Mal:}
\begin{itemize}
\item Das iPhone erkennt die Webseite/App
\item Bietet an, das Passwort \textbf{einzusetzen}
\item Du bestätigst mit Face ID / Touch ID
\item Fertig -- eingeloggt!
\end{itemize}

\columnbreak

\section*{\textcolor{shieldgray}{Passwörter verwalten}}

\textbf{Alle Passwörter ansehen:}\\
Einstellungen → \textbf{Passwörter}\\
(Entsperren mit Face ID / Touch ID / Code)

\vspace{2mm}

\textbf{Was du dort findest:}
\begin{itemize}
\item Liste aller gespeicherten Passwörter
\item Webseite/App + Benutzername + Passwort
\item \textbf{Sicherheitsempfehlungen} (!)
\end{itemize}

\textbf{Sicherheitsempfehlungen beachten:}\\
Das iPhone warnt dich, wenn:
\begin{itemize}
\item Ein Passwort zu \textbf{schwach} ist
\item Du ein Passwort \textbf{mehrfach} verwendest
\item Ein Passwort in einem \textbf{Datenleck} aufgetaucht ist
\end{itemize}

$\rightarrow$ Diese Passwörter \textbf{sofort ändern}!

\vspace{3mm}

\begin{center}
\fcolorbox{keyblue}{keyblue!8}{%
\begin{minipage}{0.9\linewidth}
\vspace{2mm}
\textbf{Tipp: Passwort teilen}\\[2pt]
Du kannst Passwörter sicher per \textbf{AirDrop} an Familienmitglieder senden -- ohne sie laut zu sagen.
\vspace{2mm}
\end{minipage}}
\end{center}

\vspace{3mm}

\begin{center}
\fcolorbox{warnred}{warnred!8}{%
\begin{minipage}{0.9\linewidth}
\vspace{2mm}
\textbf{Wichtig: Geräte-Code schützt alles!}\\[2pt]
Wer deinen Geräte-Code kennt, kann \textbf{alle Passwörter} sehen. Halte ihn geheim!
\vspace{2mm}
\end{minipage}}
\end{center}

\vspace{3mm}

\section*{Passwort vergessen?}

\begin{enumerate}
\item Auf der Webseite \glqq Passwort vergessen\grqq{} klicken
\item Link per E-Mail erhalten
\item Neues Passwort setzen
\item iPhone speichert das neue Passwort
\end{enumerate}

\end{multicols}

\newpage

% ============================================================================
% SEITE 3: ÜBUNGEN & CHECKLISTE
% ============================================================================

\begin{center}
{\LARGE\bfseries Übungen \& Checkliste}\\[3pt]
{\normalsize Teste dein Wissen über Passwörter}
\end{center}

\vspace{1mm}
\hrule
\vspace{2mm}

\begin{multicols}{2}

\textbf{\large Teil A: Gut oder schlecht?}

Bewerte diese Passwörter: G = gut, S = schlecht

\vspace{2mm}

\renewcommand{\arraystretch}{1.5}
\begin{tabular}{@{}p{5cm}cc@{}}
& G & S \\
1. \texttt{123456} & $\square$ & $\square$ \\
2. \texttt{Kx9!mLp\#2vQw} & $\square$ & $\square$ \\
3. \texttt{passwort} & $\square$ & $\square$ \\
4. \texttt{MeinHund2015} & $\square$ & $\square$ \\
5. \texttt{IljM@7Uh!kT3} & $\square$ & $\square$ \\
6. \texttt{qwertz} & $\square$ & $\square$ \\
\end{tabular}
\renewcommand{\arraystretch}{1}

\textit{\footnotesize Lösungen: S, G, S, S, G, S}

\vspace{4mm}

\textbf{\large Teil B: Richtig oder Falsch?}

\vspace{2mm}

\renewcommand{\arraystretch}{1.4}
\begin{tabular}{@{}p{7.5cm}cc@{}}
& R & F \\
1. Man sollte überall das gleiche Passwort nutzen. & $\square$ & $\square$ \\
2. Der Schlüsselbund speichert Passwörter verschlüsselt. & $\square$ & $\square$ \\
3. Ein gutes Passwort hat mindestens 12 Zeichen. & $\square$ & $\square$ \\
4. Face ID kann Passwörter automatisch ausfüllen. & $\square$ & $\square$ \\
5. Passwörter sollte man auf einem Zettel notieren. & $\square$ & $\square$ \\
\end{tabular}
\renewcommand{\arraystretch}{1}

\textit{\footnotesize Lösungen: F, R, R, R, F}

\vspace{4mm}

\textbf{\large Teil C: Erstelle ein sicheres Passwort}

Nutze den Satz-Trick:

\vspace{2mm}

Mein Satz: \rule{6cm}{0.4pt}

\rule{7cm}{0.4pt}

\vspace{2mm}

Mein Passwort (Anfangsbuchstaben + Zahlen + Sonderzeichen):

\rule{7cm}{0.4pt}

\columnbreak

\textbf{\large Teil D: Praktische Übungen}

Führe diese Aktionen auf deinem iPhone aus:

\vspace{2mm}

\begin{itemize}[label=$\square$, itemsep=3pt]
\item Öffne \textbf{Einstellungen → Passwörter}
\item Entsperre mit Face ID / Touch ID / Code
\item Schau dir die Liste der gespeicherten Passwörter an
\item Tippe auf \textbf{Sicherheitsempfehlungen}
\item Prüfe, ob du Passwörter ändern solltest
\item Prüfe in Einstellungen → [Dein Name] → iCloud, ob \textbf{Schlüsselbund} aktiviert ist
\end{itemize}

\vspace{4mm}

\textbf{\large Checkliste: Das weiß ich jetzt}

\vspace{2mm}

\begin{center}
\fcolorbox{lockgreen}{lockgreen!10}{%
\begin{minipage}{0.95\linewidth}
\vspace{2mm}
\textbf{\textcolor{lockgreen}{Sichere Passwörter}}
\begin{itemize}[label=$\square$, itemsep=2pt]
\item Ich weiß, was ein \textbf{sicheres Passwort} ausmacht
\item Ich kenne den \textbf{Satz-Trick}
\item Ich weiß, dass jedes Konto ein \textbf{eigenes Passwort} braucht
\end{itemize}
\vspace{2mm}
\end{minipage}}
\end{center}

\vspace{2mm}

\begin{center}
\fcolorbox{keyblue}{keyblue!10}{%
\begin{minipage}{0.95\linewidth}
\vspace{2mm}
\textbf{\textcolor{keyblue}{Schlüsselbund}}
\begin{itemize}[label=$\square$, itemsep=2pt]
\item Ich weiß, was der \textbf{iCloud-Schlüsselbund} ist
\item Ich kann mir Passwörter \textbf{vorschlagen} lassen
\item Ich weiß, wo ich meine Passwörter \textbf{finde}
\item Ich beachte die \textbf{Sicherheitsempfehlungen}
\end{itemize}
\vspace{2mm}
\end{minipage}}
\end{center}

\vspace{2mm}

\begin{center}
\fcolorbox{bordergray}{white}{%
\begin{minipage}{0.95\linewidth}
\vspace{2mm}
\textbf{Meine Notizen}\\[2pt]
Schlüsselbund aktiviert? $\square$ Ja \quad $\square$ Nein\\[2mm]
Anzahl gespeicherter Passwörter: \dotfill\\[2mm]
Passwörter mit Warnung: \dotfill
\vspace{2mm}
\end{minipage}}
\end{center}

\end{multicols}

\vfill

\begin{center}
\footnotesize\textit{Stand: \today{} -- Jedes Konto ein eigenes, starkes Passwort. Der Schlüsselbund merkt sie sich für dich!}
\end{center}

\end{document}
